\section{Applications of MC}
\subsection{PageRank}
Given a series of webpages $S$, each having $k_i$ hyperlinks to other webpages,
consider a random walk on the stationary MC:
\begin{itemize}
  \item State Space $S$, Time index $\mathbb{N}$
  \item $P_{ij}=\begin{cases}\frac{1}{k_i}&\text{if there is a hyperlink }i\to
    j\\0&\text{otherwise}\end{cases}$
  \item $\pi^{\left( 0 \right) }=\left( \frac{1}{\left| S \right|
    },\ldots,\frac{1}{\left| S \right| } \right) $
\end{itemize}
If the MC is nonregular, consider instead
\[\tilde{\mathbf{P}}=\lambda\mathbf{P}+\left( 1-\lambda \right) \pi_0
\mathbf{1}^{T}\]
\[\pi_{n+1}=\lambda\pi_{n}\mathbf{P}+\left( 1-\lambda \right) \pi_0.\]
for some damping factor $\lambda\in \left( 0,1 \right) $.  The limiting
distribution $\pi$ is used to rank in order of decreasing limiting probability.
\subsection{Markov Chain Monte Carlo}
Given a target distribution $\mathbf{p}=c_{n}\mathbf{p'}$, we find a MC such
that $\pi=\mathbf{p}$, using the \textcolor{Bittersweet}{Metropolis-Hastings
Algorithm}:
\begin{enumerate}
  \item Initialise with any initial state $X_0$, and an ergodic MC with the
    same sample space as $\mathbf{p}$ and transition probability matrix
    $\mathbf{Q}$
    \[\mathbf{Q}_{ij}=P\left( Y_{n+1}=j | Y_{n}=i \right) \]
  \item Using $\mathbf{Q}_{X_{n}}$, draw a proposal state $Y$.
  \item Calculate the acceptance rate of the proposal state
    \[\alpha\left( X_{n},Y \right)=\min{\left( \frac{p'_Y}{p'_{X_{n}}} \cdot
    \frac{Q_{Y,X_{n}}}{Q_{X_{n}, Y}}, 1\right) }.\] 
  \item If the proposal state is accepted, $X_{n+1}=Y$.
\end{enumerate}
On a \textcolor{Blue}{continuous} distribution $f=c_{n}f'$, the state
space $S$ is a continuous set, the transition probability is the
\textcolor{Blue}{density function} $g_x\left( y \right) =g\left( y|x \right) $
and the \textcolor{Blue}{acceptance rate} is the analogue
\[\min{\left( \frac{f'\left( Y \right)}{f'\left( X_{n} \right)}\cdot
\frac{g_Y\left( X_{n} \right) }{g_{X_{n}}\left( Y \right)}, 1\right)}\]
For sufficiently large $M$ and $gap$, we can simulate sampling from $p$ from
the MC to simulate $p$ with two approaches.
\vspace{-4mm}
 \begin{figure}[H]
  \centering
  \includegraphics[height=0.06\textwidth]{mcmc-sampling0}
  \includegraphics[height=0.06\textwidth]{mcmc-sampling1}
\end{figure}
\vspace{-4mm}
\subsection{Branching Process}
A \textcolor{Bittersweet}{branching process} is a stochastic process ${\left\{
X_{n} \right\}}_{n=0}^{\infty}$ where, for some IID RVs $\xi_i^{(j)}$ over the
sample space $\mathbb{N}$,
\[X_{n+1}=\sum_{i=0}^{X_{n}} \xi_i^{(n)}.\]
\textcolor{Bittersweet}{Properties} of Branching Processes: A
\textcolor{Blue}{stationary MC} of known transition probability function,
\textcolor{Blue}{aperiodic} with two state classes: $\left\{ 0 \right\}$ and $
\mathbb{N}^{*}$.
\begin{gather*}
\phi_{X_{n+1}|X_{n}}\left( t \right)=\phi_{\xi}^{X_{n}}\left( t \right)\\
\phi_{X_{n+1}}\left( t \right) =E\left[ \phi_\xi^{X_{n}}\left( t \right)
\right]=\phi_{X_{n}}\left( \phi_\xi\left( t \right) \right)
\end{gather*}
\[\text{If }X_0=k, \phi_{X_{n}}\left( t \right) ={\left[ \phi_\xi^{(n)}\left( t
\right) \right]}^{k}.\]
With the known \textcolor{Blue}{moments of $\xi$},
\[E\left[ X_{n}\right] =\mu E\left[ X_{n-1} \right]=\mu^{n}.\] 
\[Var\left( X_{n} \right)=\mu^{n-1}\sigma^{2}\times
\begin{cases}\frac{1-\mu^{n}}{1-\mu}&\mu\neq 1\\n&\mu=1\end{cases}.\]
\paragraph{Extinction Probability}
For absorbing state $0$, we define the stopping time and denote
\[T=\min{\left\{ n: X_{n}=0 \right\} },\quad u_{n}^{\left( k \right) }:=P\left(
T\le n|X_0=k \right) \]
Since $\left\{ T\le n \right\} =\cap ^{k}_{i=1}\left\{ T\le n \right\} \implies
u_{n}^{(k)}=u_{n}^{k}$,
\[u_{n}=\sum_{l\in S} P\left( \xi=l \right) u_{n-1}^{l}=\phi_\xi\left( u_{n-1}
\right),\]
and $u_1=P\left( \xi=0 \right) $, we can find the
\textcolor{Bittersweet}{extinction probability before generation $n$}
$u_{n}^{\left( k \right) }$ by
\[u_{n}^{(k)}=u_{n}^{k}={\left[ \phi_\xi^{(n-1)}\left( P\left( \xi=0 \right) \right) \right]}^{k}.\]
The eventual \textcolor{Bittersweet}{extinction probability} exists as it is a
non-decreasing probability wrt. $n$. 
$\phi_\xi\left( 0 \right) =P\left( \xi=0 \right) $, $\phi_\xi\left( 1 \right)
=1$, and it is a \textcolor{Blue}{convex increasing function} wrt $x$, we have
the classes of solutions:
\begin{itemize}
  \item $P\left( \xi=0 \right) =0\implies u_{\infty}=0$
  \item $P\left( \xi=0 \right) >0$, $E\left[ \xi \right] <
    1\implies u_{\infty}=1$. \textcolor{Bittersweet}{Subcritical}.
  \item $P\left( \xi=0 \right) >0$, $E\left[ \xi \right] =
    1\implies u_{\infty}=1$. \textcolor{Bittersweet}{Critical}.
  \item $P\left( \xi=0 \right) >0$, $E\left[ \xi \right] >
    1\implies u_{\infty}<1$. \textcolor{Bittersweet}{Supercritical}.
    \[u_{\infty}\text{ is the solution to }x=\phi_\xi\left( x \right).\]
\end{itemize}

